\documentclass[12pt, a4paper]{article}

% ============================================================
% PACCHETTI
% ============================================================
\usepackage[T1]{fontenc}
\usepackage[utf8]{inputenc}
\usepackage[italian]{babel}
\usepackage{lmodern}
\usepackage{microtype}
\usepackage{geometry}
\usepackage{setspace}
\usepackage{titlesec}
\usepackage{fancyhdr}
\usepackage{hyperref}
\usepackage{csquotes}
\usepackage[style=authoryear, backend=biber, citestyle=authoryear]{biblatex}
\usepackage{booktabs}
\usepackage{enumitem}
\usepackage{xcolor}
\usepackage{tcolorbox}
\usepackage{parskip}
\usepackage{footmisc}
\usepackage{emptypage}
\usepackage{lastpage}

% ============================================================
% IMPOSTAZIONI GEOMETRIA
% ============================================================
\geometry{
  top=2.5cm,
  bottom=2.5cm,
  left=3cm,
  right=2.5cm,
  headheight=15pt
}

% ============================================================
% BIBLIOGRAFIA
% ============================================================
\addbibresource{ref.bib}

% ============================================================
% COLORI
% ============================================================
\definecolor{legaleblue}{RGB}{20,52,100}
\definecolor{legalegray}{RGB}{90,90,90}
\definecolor{legalelightblue}{RGB}{220,232,246}
\definecolor{legalered}{RGB}{150,30,30}

% ============================================================
% HYPERREF
% ============================================================
\hypersetup{
  colorlinks=true,
  linkcolor=legaleblue,
  citecolor=legaleblue,
  urlcolor=legaleblue,
  pdfauthor={},
  pdftitle={La raccolta illecita di dati tramite form online},
  pdfsubject={Informatica Forense -- Diritto e protezione dei dati personali},
  pdfkeywords={GDPR, dati personali, form online, art. 167, Codice Privacy}
}

% ============================================================
% TITOLI DELLE SEZIONI
% ============================================================
\titleformat{\section}
  {\normalfont\large\bfseries\color{legaleblue}}
  {\thesection.}{0.5em}{}[\vspace{-0.3em}\rule{\linewidth}{0.4pt}]

\titleformat{\subsection}
  {\normalfont\normalsize\bfseries\color{legaleblue}}
  {\thesubsection.}{0.5em}{}

\titleformat{\subsubsection}
  {\normalfont\normalsize\itshape\color{legalegray}}
  {\thesubsubsection.}{0.5em}{}

% ============================================================
% INTESTAZIONE E PIÈ DI PAGINA
% ============================================================
\pagestyle{fancy}
\fancyhf{}
\renewcommand{\headrulewidth}{0.4pt}
\fancyhead[L]{\small\textit{La raccolta illecita di dati tramite form online}}
\fancyhead[R]{\small\textit{Informatica Forense}}
\fancyfoot[C]{\small -- \thepage\ di \pageref{LastPage} --}

% ============================================================
% BOX PER LE NORME
% ============================================================
\tcbuselibrary{skins, breakable}

\newtcolorbox{normabox}[1][]{
  colback=legalelightblue!60,
  colframe=legaleblue,
  fonttitle=\small\bfseries,
  title={#1},
  breakable,
  left=4pt, right=4pt, top=4pt, bottom=4pt,
  boxrule=0.6pt
}

\newtcolorbox{sentenzabox}[1][]{
  colback=gray!8,
  colframe=legalegray,
  fonttitle=\small\bfseries\color{legalegray},
  title={#1},
  breakable,
  left=4pt, right=4pt, top=4pt, bottom=4pt,
  boxrule=0.6pt
}

% ============================================================
% INTERLINEA
% ============================================================
\setstretch{1.3}

% ============================================================
% DOCUMENTO
% ============================================================
\begin{document}

% ------------------------------------------------------------
% PAGINA DEL TITOLO
% ------------------------------------------------------------
\begin{titlepage}
  \centering
  \vspace*{1.5cm}

  {\color{legaleblue}\rule{\linewidth}{1.2pt}}\vspace{0.3cm}
  {\color{legaleblue}\rule{\linewidth}{0.4pt}}

  \vspace{0.8cm}
  {\Huge\bfseries\color{legaleblue}
    La raccolta illecita di dati\\[0.4em]
    tramite \emph{form} online\par}

  \vspace{0.5cm}
  {\large\itshape\color{legalegray}
    Profili giuridici e tecnici della raccolta di dati personali\\
    mediante moduli web inadeguatamente predisposti\par}

  \vspace{0.8cm}
  {\color{legaleblue}\rule{\linewidth}{0.4pt}}\vspace{0.3cm}
  {\color{legaleblue}\rule{\linewidth}{1.2pt}}

  \vfill

  {\normalsize\color{legalegray}
    \textbf{Informatica Forense}\\[1em]
    \textit{Martino Papa, Alessio Minati}}

  \vspace{1.5cm}

  {\small\color{legalegray} Anno Accademico 2025/2026}
\end{titlepage}

\thispagestyle{empty}
\newpage

% ============================================================
% SEZIONE 1 – INTRODUZIONE
% ============================================================
\section{Introduzione: la raccolta di dati nel contesto digitale}
\label{sec:intro}

L'era digitale ha radicalmente mutato le modalità con cui soggetti pubblici e
privati entrano in contatto con i propri utenti. L'utilizzo di \emph{form} online 
(ossia di moduli web strutturati per la raccolta di informazioni) è diventato
prassi ordinaria in ogni settore. Tali strumenti consentono di acquisire, con
semplicità e a costo pressoché nullo, una mole considerevole di dati personali.

Proprio questa facilità di accesso ai dati costituisce, tuttavia, il principale
fattore di rischio. L'importanza dei dati personali nell'epoca attuale risiede nel fatto che
essi veicolano la nostra immagine nella rete; il diritto alla protezione
dei dati personali si configura pertanto non soltanto come espressione del
diritto alla riservatezza, ma anche come vera e propria espressione
del \emph{diritto all'autodeterminazione informativa}
\autocite[p.~38, \S~2.7]{pietropaoli2022}.

Quando il \emph{form} online viene predisposto in assenza di un'adeguata base
giuridica, senza un'informativa chiara, senza la distinzione tra dati ordinari e
dati particolari (ovvero dati \emph{sensibili}), oppure in violazione dei principi tecnici
di sicurezza, si realizzano le condizioni per una pluralità di illeciti, tanto
amministrativi quanto penali, che saranno esaminati nel presente documento.

La disciplina di riferimento è costituita principalmente dal Regolamento
(UE) 2016/679 (\emph{General Data Protection Regulation}, d'ora in poi GDPR)
\autocite{gdpr2016} e dal Decreto Legislativo 30 giugno 2003, n.~196
(\emph{Codice in materia di protezione dei dati personali}, d'ora in poi Codice
Privacy), come modificato dal D.Lgs. 10 agosto 2018, n.~101
\autocite{codiceprivacy2003, dlgs101_2018}. A livello di diritto fondamentale,
vengono in rilievo l'art.~8 della Carta dei Diritti Fondamentali dell'Unione
Europea \autocite{carta_nizza} e l'art.~8 della Convenzione Europea per la
salvaguardia dei Diritti dell'Uomo e delle Libertà Fondamentali (CEDU)
\autocite{cedu}.

% ============================================================
% SEZIONE 2 – NORMATIVA
% ============================================================
\section{Il quadro normativo}
\label{sec:normativa}

\subsection{La distinzione tra dati personali ordinari e dati particolari (sensibili)}
\label{sub:dati}

Preliminare a ogni valutazione giuridica è la corretta qualificazione del tipo
di dato raccolto tramite il \emph{form}. L'art.~4, n.~1 del GDPR definisce \textbf{dato
personale} \textquote{qualsiasi informazione riguardante una persona fisica
identificata o identificabile}.

Il legislatore europeo distingue poi una categoria di dati che richiedono una
tutela rafforzata, denominati \textbf{dati particolari} (art.~9 GDPR) o, nel
linguaggio del Codice Privacy, \textbf{dati sensibili}. Appartengono a tale
categoria:

\begin{itemize}[leftmargin=1.5cm]
  \item dati che rivelano l'origine razziale o etnica;
  \item dati relativi alle opinioni politiche, alle convinzioni religiose o filosofiche;
  \item dati relativi all'appartenenza sindacale;
  \item dati genetici e dati biometrici (se diretti a identificare
        univocamente una persona);
  \item dati relativi alla salute;
  \item dati relativi alla vita sessuale o all'orientamento sessuale.
\end{itemize}

L'art.~9, par.~1 GDPR stabilisce un divieto generale di trattamento di tali
dati, salve le eccezioni tassativamente elencate al par.~2 (consenso esplicito,
obblighi di diritto del lavoro, interessi vitali, ecc.). Ciò implica che un
\emph{form} web che, anche solo facoltativamente, contenga campi relativi a
dati particolari senza un'esplicita base giuridica ai sensi dell'art.~9 par.~2 è da
considerarsi già, di per sé, in violazione del Regolamento.

\begin{normabox}[Art.~9 GDPR -- Trattamento di categorie particolari di dati personali (sintesi)]
  \small
  \textbf{Par. 1}: ``È vietato trattare dati personali che rivelino l'origine
  razziale o etnica, le opinioni politiche, le convinzioni religiose o filosofiche,
  o l'appartenenza sindacale, nonché trattare dati genetici, dati biometrici intesi
  ad identificare in modo univoco una persona fisica, dati relativi alla salute o
  alla vita sessuale o all'orientamento sessuale della persona.''\\[0.3em]
  \textbf{Par. 2}: Il divieto non si applica nei casi ivi tassativamente previsti,
  tra cui il \emph{consenso esplicito} dell'interessato (lett.~a).
  \newline
  {\scriptsize \textit{Fonte: Regolamento (UE) 2016/679, art.~9 \autocite{gdpr2016}.}}
\end{normabox}

\subsection{Il profilo civilistico}
\label{sub:civile}

\subsubsection{I principi del GDPR applicabili alla raccolta tramite form}
\label{subsub:principi}

L'art. 5 del GDPR definisce i principi essenziali per il trattamento dei dati, 
che devono essere rispettati sin dalla fase di creazione di qualsiasi \emph{form} online.
In particolare:

\begin{enumerate}[label=\alph*), leftmargin=1.5cm]
  \item \textbf{Liceità, correttezza e trasparenza} (art.~5, par.~1, lett.~a):
        ogni raccolta deve fondarsi su una base giuridica valida (consenso,
        contratto, obbligo legale, interessi vitali, interesse pubblico o legittimo
        interesse del titolare);

  \item \textbf{Limitazione della finalità} (art.~5, par.~1, lett.~b):
        i dati raccolti tramite il modulo possono essere trattati esclusivamente
        per le finalità determinate, esplicite e legittime indicate nell'informativa;

  \item \textbf{Minimizzazione dei dati} (art.~5, par.~1, lett.~c):
        il \emph{form} deve raccogliere soltanto i dati \emph{adeguati, pertinenti
        e limitati} a quanto necessario rispetto alle finalità perseguite;
        qualsiasi campo superfluo integra una violazione del Regolamento;

  \item \textbf{Esattezza} (art.~5, par.~1, lett.~d): il titolare deve
        adottare misure ragionevoli per garantire la correttezza dei dati
        raccolti;

  \item \textbf{Limitazione della conservazione} (art.~5, par.~1, lett.~e):
        i dati devono essere conservati in una forma che consenta l'identificazione
        degli interessati per un periodo non superiore al conseguimento delle
        finalità per le quali sono trattati;

  \item \textbf{Integrità e riservatezza} (art.~5, par.~1, lett.~f):
        il \emph{form} deve essere tecnicamente protetto, con connessione
        cifrata (protocollo HTTPS/TLS) e misure di sicurezza adeguate.
\end{enumerate}

\subsubsection{L'obbligo di informativa e il consenso}
\label{subsub:consenso}

L'art.~13 GDPR impone al titolare del trattamento di fornire all'interessato,
al momento della raccolta dei dati, una serie di informazioni in forma
concisa, trasparente e intellegibile. Un \emph{form} online sprovvisto di
informativa, oppure che rimandi a un testo generico e incomprensibile, è in
violazione di tale obbligo.

Qualora la raccolta dei dati si basi sul \textbf{consenso} dell'utente, non essendo strettamente 
obbligatoria per altri motivi (art.~6, par.~1, lett.~a), 
l'art.~7 del GDPR stabilisce che tale consenso debba essere:

\begin{itemize}[leftmargin=1.5cm]
  \item \textbf{libero}: l'accesso al servizio non deve essere condizionato alla
        prestazione del consenso per finalità non necessarie;
  \item \textbf{specifico}: un unico checkbox che raggruppi più finalità non
        soddisfa il requisito;
  \item \textbf{informato}: deve riferirsi a un'informativa adeguata;
  \item \textbf{inequivocabile}: le caselle pre-spuntate (\emph{pre-ticked boxes})
        non integrano un consenso valido.
\end{itemize}

\subsubsection{Responsabilità civile del titolare del trattamento}
\label{subsub:responsabilitacivile}

L'art.~82 GDPR sancisce il diritto di ogni interessato che abbia subito
un danno --- materiale o immateriale --- in conseguenza di una violazione del
Regolamento, di ottenere il risarcimento dal titolare o dal responsabile del
trattamento. Il titolare è esonerato da responsabilità soltanto se dimostri
che l'evento dannoso non gli è in alcun modo imputabile: si tratta di una
\emph{presunzione di responsabilità} con onere della prova invertito a carico
del titolare.

Le sanzioni amministrative pecuniarie, irrogate dal Garante per la protezione
dei dati personali, possono raggiungere, ai sensi dell'art.~83 GDPR:
\begin{itemize}[leftmargin=1.5cm]
  \item fino a \textbf{10 milioni di euro} o, per le imprese, il \textbf{2\%}
        del fatturato mondiale annuo, per le violazioni di cui al par.~4
        (es.\ violazione delle condizioni del consenso, mancata informativa);
  \item fino a \textbf{20 milioni di euro} o il \textbf{4\%} del fatturato
        mondiale annuo, per le violazioni di cui al par.~5 (es.\ violazione
        dei principi di base, trattamento di dati particolari senza base giuridica).
\end{itemize}

\subsection{Il profilo penalistico}
\label{sub:penale}

\subsubsection{Il trattamento illecito di dati personali (art.~167 Codice Privacy)}
\label{subsub:art167}

Per fronteggiare le violazioni della sfera
della riservatezza nell'ambito digitale, nel nostro ordinamento sono
state introdotte alcune norme penali, inserite non nel Codice penale ma nel
d.lgs. 196/2003 (Codice privacy)
\autocite[p.~38, \S~2.7]{pietropaoli2022}.

L'art.~167 del Codice Privacy, nella formulazione successiva all'adeguamento
al GDPR operato dal D.Lgs.~101/2018, punisce con la reclusione da
\textbf{sei mesi a un anno e sei mesi} chiunque, al fine di trarne per sé o
per altri profitto ovvero al fine di recare ad altri un danno, proceda al
trattamento di dati personali in violazione delle disposizioni di cui agli
artt.~123, 126, 129 e 130 del medesimo decreto, nonché in violazione
delle misure di garanzia di cui all'art.~2-\emph{quinquiesdecies}, laddove ne consegua una lesione.

La pena è aumentata a \textbf{uno a tre anni} qualora il fatto riguardi
dati sensibili.

\begin{normabox}[Art.~167 D.Lgs. 196/2003 -- Trattamento illecito di dati (sintesi)]
  \small
  ``1. Salvo che il fatto costituisca più grave reato, chiunque, al fine di
  trarne per sé o per altri profitto ovvero al fine di recare ad altri un danno,
  procede al trattamento di dati personali in violazione di quanto disposto
  \textelp{} è punito, se dal fatto deriva nocumento, con la reclusione da sei
  mesi a un anno e sei mesi. \\
  2. Se il fatto di cui al comma 1 riguarda dati
  genetici, biometrici o relativi alla salute, alla vita sessuale, all'orientamento
  sessuale o alle opinioni, \textelp{} la pena è la reclusione da uno a tre anni.''
  \newline
  {\scriptsize \textit{Fonte: D.Lgs. 30 giugno 2003, n.~196, art.~167
  \autocite{codiceprivacy2003}.}}
\end{normabox}

Perché si configuri questo reato, la giurisprudenza ha stabilito che non è strettamente necessario 
un danno economico: è sufficiente anche un danno di natura morale. Inoltre, il vantaggio illecito 
perseguito dal colpevole può consistere in un qualsiasi tipo di beneficio, anche non monetario.

\subsubsection{La comunicazione e diffusione illecita di dati personali su larga scala (art.~167-bis)}
\label{subsub:art167bis}

L'art.~167-\emph{bis} del Codice Privacy (introdotto dal D.Lgs.~101/2018)
punisce con la reclusione da \textbf{uno a sei anni} chiunque comunichi o
diffonda dati personali oggetto di trattamento su larga scala, in violazione
delle disposizioni del GDPR. La norma si applica con particolare intensità
nel caso in cui i dati raccolti tramite \emph{form} vengano ceduti a terzi o
pubblicati senza autorizzazione, configurando un cosidetto \textbf{data breach
deliberato}.

\begin{normabox}[Art.~167-bis D.Lgs. 196/2003 -- Comunicazione e diffusione illecita di dati su larga scala (sintesi)]
  \small
  ``1. Salvo che il fatto costituisca più grave reato, chiunque comunica o diffonde,
  al fine di trarne profitto per sé o altri ovvero al fine di arrecare danno, un
  archivio automatizzato o una parte sostanziale di esso contenente dati personali
  oggetto di trattamento su larga scala, in violazione degli articoli 2-ter, 2-sexies
  e 2-octies, è punito con la reclusione da uno a sei anni.''
  \newline
  {\scriptsize \textit{Fonte: D.Lgs. 30 giugno 2003, n.~196, art.~167-bis
  \autocite{codiceprivacy2003}.}}
\end{normabox}

\subsubsection{L'acquisizione fraudolenta di dati personali su larga scala (art.~167-ter)}
\label{subsub:art167ter}

L'art.~167-\emph{ter} del Codice Privacy punisce con la reclusione da
\textbf{uno a quattro anni} chiunque, al fine di trarne profitto per sé
o per altri ovvero per recare danno, acquisisca fraudolentemente dati
personali oggetto di trattamento su larga scala. Tale fattispecie riguarda,
tra l'altro, la predisposizione di \emph{form} ingannevoli (c.d.\ \emph{dark
patterns}) che inducano l'utente a fornire dati personali in modo che non è
consapevole o che non ha inteso.

\begin{normabox}[Art.~167-ter D.Lgs. 196/2003 -- Acquisizione fraudolenta di dati su larga scala (sintesi)]
  \small
  ``1. Salvo che il fatto costituisca più grave reato, chiunque, al fine di trarne
  profitto per sé o altri ovvero di arrecare danno, acquisisca con mezzi
  fraudolenti un archivio automatizzato o una parte sostanziale di esso contenente
  dati personali oggetto di trattamento su larga scala è punito con la reclusione
  da uno a quattro anni.''
  \newline
  {\scriptsize \textit{Fonte: D.Lgs. 30 giugno 2003, n.~196, art.~167-ter
  \autocite{codiceprivacy2003}.}}
\end{normabox}

% ============================================================
% SEZIONE 3 – CASI PRATICI
% ============================================================
\section{Casi pratici e giurisprudenza}
\label{sec:casipratici}

\subsection{Garante Privacy -- Sanzione per form web senza adeguata informativa (2021)}
\label{sub:garante2021}

\begin{sentenzabox}[Garante per la protezione dei dati personali, Provv. 29 aprile 2021]
  Un ente pubblico aveva predisposto sul proprio sito web un modulo di
  contatto (\emph{form}) che raccoglieva nome, cognome, indirizzo e-mail
  e numero di telefono degli utenti, senza che nella pagina fosse presente
  alcuna informativa adeguata ai sensi dell'art.~13 GDPR. Il Garante ha
  accertato la violazione dei principi di liceità, correttezza e trasparenza
  (art.~5, par.~1, lett.~a GDPR) e ha irrogato una sanzione pecuniaria
  amministrativa, disponendo altresì la rettifica del \emph{form}.
  \newline
  {\scriptsize \textit{Fonte: \autocite{garante2021_newsletter}.}}
\end{sentenzabox}

\medskip

Il caso illustra la fattispecie più frequente: il titolare
del trattamento non predispone alcuna informativa visibile nella pagina che
ospita il modulo, o la inserisce in un link oscuro nel \emph{footer} del sito,
ritenendo erroneamente di assolvere così all'obbligo informativo. Tale prassi
non soddisfa il requisito della \emph{trasparenza} sancito dall'art.~12 GDPR,
il quale impone che le informazioni siano fornite in forma concisa, facilmente
accessibile e comprensibile.

\subsection{Garante Privacy -- Trattamento illecito di dati particolari tramite form (2022)}
\label{sub:garante2022}

\begin{sentenzabox}[Garante per la protezione dei dati personali, Provv. 15 settembre 2022]
  Un'azienda privata aveva inserito nel proprio \emph{form} di candidatura
  online campi contrassegnati come \emph{facoltativi} relativi allo stato di
  salute e all'eventuale appartenenza sindacale dei candidati. Il Garante ha
  rilevato che, ai sensi dell'art.~9 GDPR, il trattamento di dati particolari
  non può essere fondato esclusivamente sulla facoltatività del campo, ma
  richiede un'esplicita e valida base giuridica tra quelle elencate all'art.~9,
  par.~2. L'ente è stato sanzionato per violazione degli artt.~5, 6 e 9 GDPR.
  \newline
  {\scriptsize \textit{Fonte: \autocite{garante2022_form}.}}
\end{sentenzabox}

\medskip

Tale pronuncia evidenzia un equivoco diffuso: la \emph{facoltatività} di un
campo non esclude la necessità di una base giuridica per il trattamento dei
dati ivi inseriti. Se il campo è facoltativo ma raccoglie dati particolari, occorre
comunque che il consenso sia esplicito (art.~9, par.~2, lett.~a GDPR),
granulare e separato dagli altri consensi.

\subsection{Corte di Cassazione penale -- Trattamento illecito di dati a fini di profitto (2019)}
\label{sub:cass2019}

\begin{sentenzabox}[Cass. pen., sez. II, 15 gennaio 2019, n.~1285]
  La Suprema Corte ha confermato la configurabilità del reato di cui all'art.~167
  del Codice Privacy in capo a un operatore che aveva sistematicamente raccolto,
  tramite un servizio web, dati personali degli utenti senza il loro consenso e li
  aveva successivamente ceduti a terzi a fini commerciali. La Corte ha precisato
  che il \emph{fine di profitto} richiesto dalla norma è integrato anche dalla sola
  prospettiva di un vantaggio economico indiretto, e che il \emph{nocumento}
  può consistere anche in un danno non patrimoniale, quale la perdita del controllo
  sui propri dati personali.
  \newline
  {\scriptsize \textit{Fonte: \autocite{cass_pen_2019_1285}.}}
\end{sentenzabox}

\medskip

Questa pronuncia è particolarmente significativa per i moduli web destinati alla profilazione 
commerciale. Essa evidenzia come la raccolta di dati in assenza di un'idonea informativa e 
di un valido consenso, seguita dalla loro cessione a terzi, sia sufficiente a configurare 
il reato di trattamento illecito, a prescindere dal guadagno effettivamente ottenuto.

\subsection{Corte di Cassazione civile -- Risarcimento del danno per trattamento illecito (2023)}
\label{sub:cass2023}

\begin{sentenzabox}[Cass. civ., sez. I, 10 gennaio 2023, n.~407]
  La Suprema Corte ha confermato la condanna al risarcimento del danno non
  patrimoniale in favore di un utente i cui dati personali erano stati raccolti
  tramite un modulo web privo di adeguata informativa e, successivamente,
  comunicati a terzi non autorizzati. I giudici hanno richiamato l'art.~82 GDPR
  e ribadito che spetta al titolare del trattamento fornire la prova liberatoria
  dell'assenza di colpa. Il mero fatto che i dati siano stati trattati in violazione
  del GDPR costituisce presupposto sufficiente per riconoscere il danno
  non patrimoniale, senza che l'interessato debba dimostrare un pregiudizio
  ulteriore concreto.
  \newline
  {\scriptsize \textit{Fonte: \autocite{cass2023_privacyform}.}}
\end{sentenzabox}

\medskip

La decisione riveste un'importanza fondamentale sul piano della responsabilità civile del titolare 
del trattamento. L'inversione dell'onere della prova e il riconoscimento del danno non patrimoniale 
legato alla semplice perdita di controllo sui propri dati personali rendono l'azione risarcitoria 
un rischio concreto per chiunque gestisca moduli web non conformi. 
La conformità del \emph{form} non rappresenta quindi solo una tutela contro le sanzioni del Garante, 
ma uno scudo indispensabile contro le richieste di risarcimento dirette degli utenti.

% ============================================================
% SEZIONE 4 – CONCLUSIONI
% ============================================================
\section{Conclusioni}
\label{sec:conclusioni}

L'analisi svolta consente di formulare le seguenti considerazioni conclusive.

In primo luogo, la predisposizione di un \emph{form} online per la raccolta di
dati personali non è un'operazione neutra sul piano giuridico, ma implica
l'assunzione di precisi obblighi di legge in capo al titolare del trattamento,
la cui violazione espone a conseguenze che spaziano dalla sanzione amministrativa
pecuniaria fino alla responsabilità penale.

In secondo luogo, la distinzione tra dati personali ordinari e dati particolari
(sensibili) è di importanza cruciale: l'inserimento nel \emph{form}, anche solo
facoltativamente, di campi relativi alla salute, all'orientamento sessuale, alle
opinioni politiche o ad altri elementi elencati nell'art.~9 GDPR richiede
una base giuridica rafforzata e misure di sicurezza tecnicamente adeguate.

Da un punto di vista operativo, chiunque intenda predisporre un \emph{form}
online per la raccolta di dati personali deve, come \emph{minima misura}:

\begin{enumerate}[label=\arabic*., leftmargin=1.5cm]
  \item definire e documentare le finalità del trattamento;
  \item identificare la base giuridica applicabile;
  \item predisporre e rendere visibile un'informativa completa ex art.~13 GDPR;
  \item raccogliere soltanto i dati strettamente necessari (\emph{data minimisation});
  \item adottare misure tecniche di sicurezza adeguate (cifratura, HTTPS, ecc.);
  \item in caso di dati particolari, acquisire un consenso esplicito e separato.
\end{enumerate}

La mancata osservanza anche di una sola di queste prescrizioni può
integrare, a seconda dei casi, un illecito amministrativo o un illecito penale,
con le gravi conseguenze patrimoniali e reputazionali che ne derivano.

% ============================================================
% BIBLIOGRAFIA
% ============================================================
\newpage
\phantomsection
\addcontentsline{toc}{section}{Riferimenti normativi e bibliografici}
\printbibliography[title={Riferimenti normativi e bibliografici}]

\end{document}
